\documentclass[11pt,a4paper]{article}
\usepackage[utf8]{inputenc}
\usepackage[T1]{fontenc}
\usepackage{lmodern}
\usepackage[margin=2.5cm]{geometry}
\usepackage{hyperref}
\usepackage{graphicx}
\usepackage{booktabs}
\usepackage{multirow}
\usepackage{xcolor}
\usepackage{listings}
\usepackage{enumitem}

\lstdefinestyle{bashstyle}{
  language=bash,
  basicstyle=\ttfamily\small,
  breaklines=true,
  frame=single,
  backgroundcolor=\color{gray!5},
  rulecolor=\color{gray!50},
  columns=fullflexible
}
\lstdefinestyle{pythonstyle}{
  language=Python,
  basicstyle=\ttfamily\small,
  breaklines=true,
  frame=single,
  backgroundcolor=\color{gray!5},
  rulecolor=\color{gray!50},
  columns=fullflexible
}

\title{PathX - HW1 Report}
\author{Juliette Tonneau \quad Minseok Kim \quad Bita Khashechian \quad Danil Smirnov\\[4pt]
\textbf{Graph Databases — Gianmaria Silvello}}

\begin{document}

\maketitle

\section{Project Description}

The goal is to build a \textbf{system that integrates an SQL database} with \textbf{a Large Language Model (LLM)}, enabling natural language query processing and query optimization inspired by the \emph{Galois} paper presented in the course.


\section{Overview and Key Insights from the GALOIS Paper}

The SIGMOD 2025 paper \emph{``Logical and Physical Optimizations for SQL Query Execution over Large Language Models''} presents \textbf{Galois}, a system that enables SQL query execution directly over Large Language Models (LLMs). The observation is that, although LLMs are capable of answering declarative queries and extracting structured data from unstructured sources, their performance declines with complex queries, often resulting in incomplete or inaccurate outputs. Existing database optimization techniques target first cost or latency reduction without adequately addressing the challenge of ensuring result quality.

Galois implements an \textbf{intermediate optimization layer} between SQL or NL queries and LLMs while treating the LLM as a \textbf{storage layer}. The system modifies both logical and physical query plans to adapt to the distinct operational characteristics of language models.

\subsection*{Key innovations we retain}

\begin{itemize}[nosep]
  \item Redefine the query optimization process to \textbf{balance execution cost and result quality}.
  \item Introduce two physical operators for retrieving data from LLMs: \emph{table-scan} and \emph{key-scan}.
  \item Dynamically build metadata and estimate confidence scores since LLMs lack catalogs.
  \item Treat the \textbf{LLM as a data source}, not only a reasoning tool.
  \item Maintain a \textbf{cost--quality} balance.
  \item Structure results in \textbf{JSON} to ensure compatibility with analysis tools.
\end{itemize}

\section{Detailed Implementation Schedule}

\subsection*{Phase 1 – Tech Selection and System Setup}
Set up the full system (architecture, Python environment, DuckDB databases) with setup tests and documentation.

\subsection*{Phase 2 – NL and SQL Baselines}
Build NL and SQL baselines, integrate into LangChain, test prompting strategies, and set up evaluation metrics (accuracy, F1).

\subsection*{Phase 3 – Galois Features (Part I)}
Implement core Galois features (LLMScan, Filter-LLMScan, joins, two-phase retrieval) and prepare a demo comparing naïve vs.\ optimized runs.

\subsection*{Phase 4 – Galois Features (Part II)}
Develop and test optimized versions (Selective, All-Attribute, Full), replicate key results, and finalize the report.


\section{Stage 1 --- Tech Selection and System Setup}

This stage focuses on:
\begin{itemize}[nosep]
  \item selecting and setting up the development environment,
  \item preparing a local SQL database (DuckDB),
  \item implementing a Python pipeline that:
    \begin{enumerate}[nosep]
      \item connects to the database,
      \item runs SQL queries from a \texttt{.sql} file,
      \item saves results in a \texttt{.json} output file in \emph{Full} format (rows + metrics) defined by the tutor.
      \item assesses our system’s outputs by comparing them with the ground-truth data.
    \end{enumerate}
\end{itemize}

\section{Technology Stack}

\begin{table}[h]
  \centering
  \begin{tabular}{@{}llp{7cm}@{}}
    \toprule
    \textbf{Component} & \textbf{Tool / Library} & \textbf{Description} \\
    \midrule
    Language & Python 3.11+ & Main programming language \\
    Database & DuckDB & Lightweight, file-based SQL DB optimized for analytics \\
    LLM (later) & IBM WatsonX.ai & LLM integration for later stages \\
    Libraries & \texttt{duckdb}, \texttt{pandas}, \texttt{sqlglot}, \texttt{python-dotenv} & Core Python dependencies \\
    IDE & Visual Studio Code & Main development environment \\
    \bottomrule
  \end{tabular}
  \caption{Technology stack.}
\end{table}


\section{Tools and Libraries}
\begin{itemize}
  \item \textbf{numpy} -- numerical support for pandas / data processing.
  \item \textbf{pandas} -- normalize LLM JSON output into tabular form before loading into DuckDB or analyzing.
  \item \textbf{duckdb} -- execute post-LLM SQL (joins, aggregations) and store intermediate results.
  \item \textbf{sqlglot} -- parse SQL, rewrite/decompose queries, and produce logical plans that map to LLM operators.
  \item \textbf{IBM watsonx.ai SDK} (\texttt{ibm-watsonx-ai}, \texttt{ibm-watsonx-ai.foundation\_models}) -- call the IBM LLM directly and implement the two operators (LLM table-scan, LLM key-scan).
  \item \textbf{pydantic} -- validate and enforce the schema of LLM JSON outputs.
  \item \textbf{python-dotenv} -- to keep IBM API keys, project IDs, deployments, etc. out of the code.
  \item \textbf{pyyaml} -- To store things like: which model to use, what prompt template to use for table-scan/key-scan, cost/quality parameters.
  \item \textbf{logging} -- record prompts, LLM latency, token usage, and validation errors for debugging and analysis.
\end{itemize}

\section{ Directories}

\begin{description}

  \item[\texttt{/data}]%
Stores datasets (e.g. .csv files), query files (.sql), and text inputs (.txt, for NL), as well as ingest.sql for each dataset. 
  
  \item[\texttt{/ground}]%
Reference data used for evaluation.

  \item[\texttt{/db}]%
Contains all the databases for datasets and it is in .gitignore.

  \item[\texttt{/scripts}]%
Files provided by tutor for evaluation and etc.

\item[\texttt{/src}]%
Project’s core codebase, divided into the modules below:
  \begin{itemize}
    \item \texttt{/db} — For initializing the DuckDB connection and ingesting CSV files into DuckDB.
    \item \texttt{/llm} — For prompting and connecting to LLM.
    \item \texttt{/pipelines} — For running queries in json, implementing EXPLAIN and saving the result in .txt format for human inspection and .json for machine processing. Also initial baselines for future stages.
    \item \texttt{/utils} — Utilities such as loggers to make the program’s behavior observable and timer for compare runs and report the evaluation.
  \end{itemize}

\item[\texttt{/submissions}]%
Final result of our system in both formats including \texttt{.json} and \texttt{.txt}, and also EXPLAIN output for displaying physical and logical plan.

  \item[\texttt{/config}]%
  Contains configuration setup files that contains:
  \begin{itemize}
    \item \texttt{.env} — Environment-specific variables (e.g. IBM watsonx credentials, API keys, service URLs).
    \item \texttt{config.yaml} — Main configuration file.
  \end{itemize}

 \item[\texttt{/test}]%
Holds the test scripts, for both small parts and full flows.

\end{description}

\section{Setup and Implementation}

\begin{enumerate}
    \item \textbf{Environment setup:} Creates a Python virtual environment and installs all dependencies from \texttt{requirements.txt}.
    \item \textbf{Data ingestion:} Loads all datasets into the database.
    \item \textbf{Query processing:} Executes all predefined queries.
    \item \textbf{Explain export:} Generates and exports EXPLAIN results in JSON format.
    \item \textbf{Evaluation:} Compares system outputs with the ground truth (provided by tutor) to verify accuracy.
\end{enumerate}

The \texttt{main()} function calls all these steps in order, providing a project initialization and validation pipeline.


\end{document}